
\chapter{Méthodologie et notions DevSecOps}
\label{ch:methodo}

\section{Introduction}

Ce chapitre fixe les concepts utilisés pour~\textbf{lire et maintenir l'application sous~\texttt{paas/frontend}} ainsi que la cadence suivie lors de son développement. Il complète le chap.~\emph{«~Planification et architecture~\ref{ch:archi}~»}.

\section*{Axes fonctionnels du livrable}

Le~\texttt{paas} concentre la valeur sur trois axes~:

\begin{enumerate}[leftmargin=*]
\item \textbf{Orchestration unifiée}~: une API~/~IHM commune pour plusieurs moteurs de build et des états de livraison homogènes~;
\item \textbf{Sécurité en chaîne}~: exposition centralisée des points de rattache~(scanners, registre, signatures, politiques)~;
\item \textbf{Traçabilité}~: suivi projet par projet jusqu'au manifeste~(commit GitOps, synchronisation Argo CD).
\end{enumerate}

\section{Concepts utilisés dans le projet}

\subsection{CI, CD et dimension sécurité}

L'\textbf{intégration continue}~(CI)~automatise build et tests après changements dans le gestionnaire de versions. La \textbf{livraison continue}~(CD)~prolonge jusqu'aux environnements cibles lorsque les qualités métier~(tests, validations)~ sont satisfaites. La \textbf{DevSecOps} ajoute ces exigences dès les premiers \emph{commits}~(revue automatique des risques, dépendances, images)~sans déporter la conformité après coup seule.

\subsection{Conteneurs et Kubernetes}

Les \textbf{images OCI~/~Docker}, stockées après build dans Harbor, sont le vecteur portable entre registre et cluster. Kubernetes orchestre déploiement, mise à jour et exposition~(Ingress)~; le planificateur équilibre pods et quotas.

\subsection{GitOps~(mode pull)~}

Une fois l'image et les manifestes désirés versionnés, \textbf{Argo CD} rapproche~(pull)~ l'état du cluster de celui décrit dans Git~(aligné sur~\ref{ch:archi})~— approche différente d'un~\emph{déploiement push} où le pipeline~(CI)~pousserait directement~\texttt{kubectl}.

\subsection{Séquence typique utilisée dans l'outil}

\textbf{Webhook~/~action utilisateur~$\rightarrow$~API~(Next)~$\rightarrow$~\emph{BuildPlanner}~$\rightarrow$~moteur CI~(Jenkins ou Tekton)~$\rightarrow$~Harbor~$\rightarrow$~mise à jour repo GitOps~$\rightarrow$~Argo CD~$\rightarrow$~Kubernetes~(cf.~\ref{ch:archi})}.

\section{Planification indicative}

Une découpe en cycles courts aide à garder la traçabilité des livraisons~(alignée sur~\ref{ch:archi})~:

\begin{enumerate}[leftmargin=*]
\item \textbf{Fondations}~: modèle projet~/~utilisateur~/~sessions, gardes~(auth) sur routes API)~;
\item \textbf{Construction}~: branchement backend CI, premières exécutions Jenkins~/~Tekton, méta-données d'exécution~;
\item \textbf{Promotion et runtime}~: GitOps, Argo CD, affichage d'étapes côté IHM)~;
\item \textbf{Sécurité et métriques}~: scanners, registre, dashboards et politiques facultatives)~.
\end{enumerate}

\section*{Méthodologie}

Le développement s'effectue dans un esprit~\textbf{itératif}~(objectifs~(backlog)~découpés, revues après chaque jalon)~, compatible avec Scrum sans lier rigidement nom d'outil à un fournisseur de gestion de tickets.

\section*{Bilan}

Ce cadre lexical et cette planification~\emph{orientée projet} retrouvent le schéma d'architecture et les artefacts logiciels~(fichiers de configuration environnementaux, backends \texttt{BUILD\_BACKEND}=\texttt{jenkins}\,|\,\texttt{tekton}, etc.) développés dans le dossier projet \texttt{paas/}.
